\begin{figure}[tb]
\centering
\begin{tikzpicture}[
  box/.style={draw,rounded corners=2pt,align=center,
    text width=3.45cm,minimum height=1.5cm,inner sep=5pt},
  arr/.style={-{Stealth[length=2mm]},line width=.65pt},
  every node/.style={font=\small}]
  \node[box] (original) at (0,0)
    {$D_n$\\tensor occurrence\\
     $(N\lambda,N\mu,N\xi)$};
  \node[box] (doubled) at (4.5,0)
    {$D_{r+1}$\\tensor occurrence\\
     $(2N\widehat\lambda,2N\widehat\mu,2N\widehat\xi)$};
  \node[box] (moment) at (9,0)
    {$D_{r+1}$\\real moments\\
     $(\widehat\lambda,\widehat\mu,\widehat\xi)$};
  \node[box] (smallpath) at (9,-3.2)
    {$D_{r+1}$\\integral schedule\\
     shape $\widehat\eta$, scale $1$};
  \node[box] (largepath) at (4.5,-3.2)
    {$D_n$\\integral schedule\\
     original shape $\eta$, scale $1$};
  \node[box] (conclusion) at (0,-3.2)
    {$D_n$\\tensor occurrence\\
     $(\lambda,\mu,\xi)$};
  \draw[arr] (original)--node[above=30pt,align=center]
    {zero-tail\\removal} (doubled);
  \draw[arr] (doubled)--node[above=30pt,align=center]
    {moments, then\\divide by $2N$} (moment);
  \draw[arr] (moment)--node[left,align=right]
    {integral-schedule\\theorem} (smallpath);
  \draw[arr] (smallpath)--node[below=30pt]
    {pad with zeros} (largepath);
  \draw[arr] (largepath)--node[below=30pt,align=center]
    {column\\intertwiners} (conclusion);
\end{tikzpicture}
\caption{The common-zero-tail branch, with $1\le r\le n-2$.
Here $\eta$ is a fundamental-column shape for the original
$D_n$ moving weight, and $\widehat\eta$ is its restriction to
the first $r+1$ coordinates. Doubling occurs only before passage
to real moments. Both integral schedules and the final tensor
occurrence use the original scale.}
\label{fig:zero-tail-scale}
\end{figure}
